\cotefile{162-5}
\scanpage{45}

$X$. De plus, on a la formule \inserted{du type $g^{\cdot}f. = f'.g'^{\cdot}$ ci-dessus, impliquant \unsure{la|le} \unsure{formule|formules}} de projection habituelle. Ceci posé, les
constructions bien connues en théorie des correspondances permettent
de définir une catégorie $\underline{B}$ dont l\unsure{\deleted{d}e|\deleted{e}d}s objets sont ceux de $C$, et les
Hom étant donnés par
\[
\mathrm{Hom}_{B}(X,Y) = B(XxY) \quad .
\]
On a un isomorphisme canonique $\underline{B} \simeq \underline{B}^{o}$, qui est l'identité sur les
objets et la symmétrie $B(XxY)\simeq B(YxX)$ sur les flêches, et un foncteur
évident \deleted{$\unsure{B|C}\quad\underline{B}$} $F. : C \to \underline{B}$, qui est l'identité sur les objets et donné
par les graphes sur les flêches. Composé avec la symmétrie précédente,
on trouve un foncteur $F^{\cdot}$: \deleted{$B$\unsure{(}} $C \relbar \underline{B}$. On vérifie enfin que le couple
$(F.,F^{\cdot})$ satisfait aux conditions envisagées dans l'alinéa 1 (ceci aussi
devrait être formalisé au cas des correspondances générales ...). \deleted{\unsure{I}}O\deleted{\unsure{l}}n
vérifie enfin qu'il a la propriété universelle qu'il faut. Lorsqu'on
se donne un couple $(F'.,F'^{\cdot})$ à val eursdans une catégorie $B'$, il faut
trouver une factorisation par $B$. Elle est toute trouvée sur les
objets, il faut la définir sur les flêches, donc définir, pour $X,Y$
dans $C$, une application
\[
\begin{array}{ll}
(*) & \qquad B(XxY) \to \mathrm{Hom}(X',Y') \quad .
\end{array}
\]
Mais si $Z$ est au\deleted{d}dessus de $XxY$, d'où $p:Z\to X$ et $q:Z\to Y$, on lui associe
un élément du deuxième membre de (*) comme le composé
\[
X' \overset{F'^{\cdot}(p)}{\longrightarrow} Z' \overset{F'.(\unsure{\deleted{p}q|\deleted{q}p})}{\longrightarrow} Y' \quad .
\]
Il faut prouver que l'élément obtenu en remp\unsure{\deleted{l}a|\deleted{a}l}çant $Z$ \inserted{$=Z_0$} par un $Z_1$ cobor\deleted{d}
dant au dessus de $X$ \inserted{est le même}. Or si $g:T\to XxY$ réalise le cobordisme, on
voit (regardant seulement $T$ et la décomposition $\mathrm{bord}(T) = Z_0\amalg Z_1$)
qu'il existe une variété $\overline{T}$ sur $S^1$ (obtenue en ``doublant'' $T$) \deleted{et une}
\deleted{application $T$ \unsure{une}} transversale au dessus d'un couple de poi\unsure{\deleted{h}n|\deleted{n}h}ts $s_o,s_1$
de $S^1$, tels que $Z_i$ soit l'image inverse \inserted{dan\unsure{s} $\overline{T}$} de $s_i$ ($i=0,1$); d'autre part,

\margin{44}
